\documentclass[12pt]{jlreq}
\usepackage{amsmath, amssymb}

\newcommand{\C}{\mathbb{C}}
\newcommand{\R}{\mathbb{R}}
\newcommand{\Z}{\mathbb{Z}}
\newcommand{\Tr}{\operatorname{Tr}}
\newcommand{\Impart}{\operatorname{Im}}
\newcommand{\AF}{\operatorname{AF}}
\newcommand{\EG}{\operatorname{EG}}
\newcommand{\AX}{\operatorname{AX}}
\newcommand{\GL}{\operatorname{GL}}
\newcommand{\Hom}{\operatorname{Hom}}
\newcommand{\Ker}{\operatorname{Ker}}
\newcommand{\Id}{\operatorname{Id}}
\newcommand{\Sym}{\operatorname{Sym}}

\begin{document}

$\Delta$ を複素平面内の単位開円板，$f : \Delta \to \mathbb{C}$ を単射正則関数とする。以下の問に答えよ。

\begin{enumerate}
  \item 正則関数 $g : \Delta \to \mathbb{C}$ で $g^2 = f'$ を満たすものが存在することを示せ。
  \item 実数 $s \in (0,1)$ に対して曲線 $\gamma_s : [0,2\pi] \to \mathbb{C}$ を
  \[
    \gamma_s(t) = f(se^{it})
  \]
  により定義する。$\gamma_s$ の長さ $L(s)$ をテイラー展開 $g(z) = \sum_{n=0}^\infty a_n z^n$ の係数 $\{a_n\}$ および $s$ を用いて表せ。
  \item $L(s)$ は $s$ の狭義単調増加関数であることを示せ。
  \item $h$ を $\Delta$ 上の正則関数とする。$0 < \delta < 1$ を満たす $\delta$ が存在して，$0 < r < \delta$ を満たす任意の $r$ に対して
  \[
    |h(0)| = \frac{1}{2\pi} \int_0^{2\pi} |h(re^{i\theta})| \, d\theta
  \]
  が成り立つとする。このとき $h$ は定数であることを示せ。
\end{enumerate}

\end{document}